% Options for packages loaded elsewhere
\PassOptionsToPackage{unicode}{hyperref}
\PassOptionsToPackage{hyphens}{url}
\documentclass[
  13pt,
]{article}
\usepackage[table]{xcolor}
\definecolor{headerblue}{RGB}{41, 65, 122}
\definecolor{rowgray}{RGB}{240, 243, 250}
\renewcommand{\arraystretch}{1.4}
\usepackage[margin=2.5cm]{geometry}
\usepackage{amsmath,amssymb}
\setcounter{secnumdepth}{-\maxdimen} % remove section numbering
\usepackage[T5]{fontenc}
\usepackage[utf8]{inputenc}
\usepackage{textcomp}
\usepackage{times}
% Use upquote if available, for straight quotes in verbatim environments
\IfFileExists{upquote.sty}{\usepackage{upquote}}{}
\IfFileExists{microtype.sty}{% use microtype if available
  \usepackage[]{microtype}
  \UseMicrotypeSet[protrusion]{basicmath} % disable protrusion for tt fonts
}{}
\makeatletter
\@ifundefined{KOMAClassName}{% if non-KOMA class
  \IfFileExists{parskip.sty}{%
    \usepackage{parskip}
  }{% else
    \setlength{\parindent}{0pt}
    \setlength{\parskip}{6pt plus 2pt minus 1pt}}
}{% if KOMA class
  \KOMAoptions{parskip=half}}
\makeatother
\usepackage{longtable,booktabs,array}
\usepackage{calc} % for calculating minipage widths
% Correct order of tables after \paragraph or \subparagraph
\usepackage{etoolbox}
\makeatletter
\patchcmd\longtable{\par}{\if@noskipsec\mbox{}\fi\par}{}{}
\makeatother
% Allow footnotes in longtable head/foot
\IfFileExists{footnotehyper.sty}{\usepackage{footnotehyper}}{\usepackage{footnote}}
\makesavenoteenv{longtable}
\usepackage[vietnamese]{babel}
% get rid of language-specific shorthands (see #6817):
\let\LanguageShortHands\languageshorthands
\def\languageshorthands#1{}
\setlength{\emergencystretch}{3em} % prevent overfull lines
\providecommand{\tightlist}{%
  \setlength{\itemsep}{0pt}\setlength{\parskip}{0pt}}
\usepackage{bookmark}
\IfFileExists{xurl.sty}{\usepackage{xurl}}{} % add URL line breaks if available
\urlstyle{same}
\hypersetup{
  pdflang={vi},
  hidelinks,
  pdfcreator={LaTeX via pandoc}}

\author{}
\date{}

\begin{document}

\textbf{BÁO CÁO TIẾN ĐỘ BUỔI 7 - NHÓM 10}

\textbf{Đề tài:} Xây dựng hệ thống nhúng hỗ trợ bán hàng thông minh

\subsection{\texorpdfstring{\textbf{I. Tóm tắt tiến độ hiện
tại}}{I. Tóm tắt tiến độ hiện tại}}\label{i.-tuxf3m-tux1eaft-tiux1ebfn-ux111ux1ed9-hiux1ec7n-tux1ea1i}

Dự án đã hoàn thành việc xây dựng luồng giao tiếp cơ bản giữa thiết bị
nhúng và máy chủ, đồng thời tích hợp thành công các mô hình AI cốt lõi.
Mô hình khử nhiễu tự phát triển cho kết quả khả quan. Hệ thống demo (bao
gồm web streaming và xử lý giọng nói) hiện đang hoạt động ổn định và
trơn tru.

\subsection{\texorpdfstring{\textbf{II. Các công việc đã hoàn
thành}}{II. Các công việc đã hoàn thành}}\label{ii.-cuxe1c-cuxf4ng-viux1ec7c-ux111uxe3-houxe0n-thuxe0nh}

\textbf{2.1. Phát triển thiết bị nhúng (Hardware/IoT)}

\begin{itemize}
\item
  Thiết lập và cấu hình ESP32 (Sử dụng ESP32-CAM và Microphone).
\item
  Tối ưu hóa tài nguyên phần cứng bằng cách chia tải xử lý trên 2 nhân
  (dual-core) của ESP32:

  \begin{itemize}
  \item
    1 core chuyên trách xử lý tín hiệu âm thanh (audio) từ mic.
  \item
    1 core chuyên trách xử lý hình ảnh (video streaming) từ camera.
  \end{itemize}
\end{itemize}

\textbf{2.2. Xây dựng Server và tích hợp các AI Services}

\begin{itemize}
\item
  Phát triển máy chủ trung tâm sử dụng \ldots\ldots.
\item
  Thiết lập và liên kết 3 services chính trên server:

  \begin{itemize}
  \item
    \textbf{Service 1 (Denoise):} Khử nhiễu âm thanh đầu vào.
  \item
    \textbf{Service 2 (STT - Speech to Text):} Chuyển đổi giọng nói
    thành văn bản sử dụng mô hình Whisper.
  \item
    \textbf{Service 3 (Classification):} Phân loại ý định/nội dung dựa
    trên văn bản sử dụng mô hình RoBERTa transformer.
  \end{itemize}
\end{itemize}

\textbf{2.3. Nghiên cứu và huấn luyện mô hình Khử nhiễu (Denoise Model)}

\begin{itemize}
\item
  \textbf{Tiền xử lý và chuẩn bị dữ liệu:} Xây dựng tập dữ liệu huấn
  luyện bằng cách sử dụng bộ dữ liệu tiếng Việt VIVOS (kích thước
  \textgreater{} 1GB) kết hợp ngẫu nhiên (mix random) với bộ dữ liệu
  nhiễu (noise) của Microsoft.
\item
  \textbf{Huấn luyện (Training):} Tự xây dựng và huấn luyện mô hình với
  dữ liệu đầu vào (input) là giọng nói có lẫn nhiễu (noisy speech) và
  mục tiêu đầu ra (output) là giọng nói sạch (clean speech).
\item
  \textbf{Đánh giá:} Thực hiện đánh giá chất lượng mô hình thông qua
  việc phân tích và so sánh phổ Fourier (Fourier spectrum) giữa tín hiệu
  nhiễu và tín hiệu sau khử nhiễu. Kết quả cho thấy mô hình xử lý khử
  nhiễu rất tốt.
\end{itemize}

\textbf{2.4. Giao diện người dùng và Demo}

\begin{itemize}
\item
  Xây dựng Web browser UI đóng vai trò là bảng điều khiển (dashboard).
\item
  Hiển thị luồng video streaming theo thời gian thực từ ESP32.
\item
  Tích hợp tính năng hiển thị trực quan dữ liệu giọng nói (Voice
  Recognition UI):

  \begin{itemize}
  \item
    Hiển thị âm thanh dưới dạng biểu đồ phổ (Spectrogram) và dạng sóng
    (Waveform).
  \item
    Cho phép quan sát và so sánh trực quan tín hiệu trước và sau khi đi
    qua mô hình khử nhiễu.
  \end{itemize}
\end{itemize}

\subsection{\texorpdfstring{\textbf{III. Kiến trúc và Cấu trúc hệ
thống}}{III. Kiến trúc và Cấu trúc hệ thống}}\label{iii.-kiux1ebfn-truxfac-vuxe0-cux1ea5u-truxfac-hux1ec7-thux1ed1ng}

\textbf{3.1. Sơ đồ khối tổng thể của hệ thống}

\textbf{3.2. Luồng hoạt động của các Services (Service Workflow)}

\textbf{3.3. Kiến trúc mô hình khử nhiễu DTLN}

\textbf{3.3.1. Tổng quan về DTLN}

DTLN (Dual-signal Transformation LSTM Network) là một mô hình học sâu
được đề xuất bởi Westhausen và Meyer (2020) nhằm giải quyết bài toán
tăng cường chất lượng giọng nói (Speech Enhancement) trong môi trường
thực tế. Mô hình được thiết kế để hoạt động hiệu quả trên các thiết bị
tài nguyên hạn chế như Raspberry Pi 4, đáp ứng yêu cầu xử lý thời gian
thực với độ trễ khoảng 32 ms.

Điểm đột phá của DTLN nằm ở kiến trúc hai giai đoạn xử lý song song
trong hai miền biểu diễn tín hiệu khác nhau. Giai đoạn đầu khai thác
miền tần số thông qua biến đổi STFT, trong khi giai đoạn hai hoạt động
trong miền đặc trưng học được (learned feature domain) thông qua bộ mã
hóa tích chập. Sự kết hợp này cho phép mô hình loại bỏ nhiễu một cách
toàn diện hơn so với các phương pháp đơn miền truyền thống.

Trong dự án này, DTLN được huấn luyện trên bộ dữ liệu clean speech từ
VIVOS (tiếng Việt) kết hợp với tập nhiễu lấy từ Microsoft Deep Noise
Suppression Challenge (DNS Challenge) của Microsoft, nhằm tối ưu hóa
hiệu năng khử nhiễu cho giọng nói tiếng Việt trong hệ thống hỗ trợ bán
hàng thông minh nhúng

\begin{table}[htbp]
\centering
\caption{Các thông số chính của mô hình DTLN}
\label{tab:dtln-params}
\begin{tabular}{>{\raggedright\arraybackslash}p{0.45\linewidth} >{\centering\arraybackslash}p{0.45\linewidth}}
\toprule
\rowcolor{headerblue} \textcolor{white}{\textbf{Thông số}} & \textcolor{white}{\textbf{Giá trị}} \\
\midrule
\rowcolor{rowgray} Tần số lấy mẫu đầu vào & 16.000 Hz, mono channel \\
Kích thước khung phân tích (blockLen) & 512 mẫu tương đương 32 ms \\
\rowcolor{rowgray} Bước dịch khung (block\_shift) & 128 mẫu tương đương 8 ms, overlap 75\% \\
Số bin tần số sau STFT & 257 bin (= blockLen / 2 + 1) \\
\rowcolor{rowgray} Số đơn vị LSTM mỗi lớp & 128 units \\
Số lớp LSTM mỗi giai đoạn & 2 lớp \\
\rowcolor{rowgray} Dropout giữa các lớp LSTM & 0,25 \\
Số bộ lọc Encoder Conv1D & 256 bộ lọc \\
\rowcolor{rowgray} Kiểu tích chập Encoder & Point-wise convolution (kernel = 1, stride = 1) \\
Hàm kích hoạt mask & Sigmoid, giá trị trong khoảng {[}0, 1{]} \\
\rowcolor{rowgray} Hàm mất mát & Negative SNR (Signal-to-Noise Ratio) \\
Phương pháp chuẩn hóa & InstantLayerNormalization (channel-wise) \\
\rowcolor{rowgray} Tổng số tham số & Khoảng 1,8 triệu tham số \\
Độ trễ xử lý & Khoảng 32 ms (real-time capable) \\
\bottomrule
\end{tabular}
\end{table}

\textbf{3.3.2. Kiến Trúc Tổng Thể}

DTLN bao gồm hai khối xử lý tuần tự, mỗi khối thực hiện ước lượng và
loại bỏ nhiễu trong một miền biểu diễn riêng biệt. Đầu vào là tín hiệu
dạng sóng (waveform) chứa giọng nói lẫn nhiễu, đầu ra là tín hiệu giọng
nói đã được phục hồi. Tín hiệu đầu ra của giai đoạn 1 được sử dụng trực
tiếp làm đầu vào cho giai đoạn 2.

Toàn bộ mô hình được tối ưu hóa \emph{end-to-end} thông qua hàm mất mát
Negative SNR bằng thuật toán tối ưu Adam với gradient clipping (clipnorm
= 3,0). Chiến lược này đảm bảo hai giai đoạn học cách phối hợp bổ sung
cho nhau mà không cần nhãn trung gian, đồng thời tránh hiện tượng bùng
nổ gradient (exploding gradient) thường gặp khi huấn luyện mạng hồi tiếp
sâu.

Hình dưới đây mô tả luồng xử lý tổng thể của mô hình, từ tín hiệu đầu
vào qua hai giai đoạn khử nhiễu đến tín hiệu đầu ra sạch.

\textbf{\emph{Hình 1. Sơ đồ kiến trúc tổng thể của mô hình DTLN}}

\textbf{x(t) → {[}STAGE 1: STFT Domain{]} → y₁(t) → {[}STAGE 2: Feature
Domain{]} → ŷ(t)}

\textbf{3.3.3. Giai Đoạn 1 - Xử Lý trong Miền Tần Số}

\textbf{3.3.3.1. Phân Khung và Biến Đổi STFT}

Tín hiệu đầu vào được chia thành các khung ngắn có kích thước 512 mẫu
(tương đương 32 ms ở tần số lấy mẫu 16 kHz) với bước dịch 128 mẫu (8
ms), tạo ra độ chồng lấn (overlap) 75\% giữa các khung liên tiếp. Biến
đổi Fourier thời gian ngắn (STFT) được áp dụng lên từng khung thông qua
phép biến đổi Fourier nửa phổ thực (rfft), thu được 257 bin tần số tương
ứng với các thành phần phổ từ 0 Hz đến 8.000 Hz.

Từ kết quả STFT, hai thành phần được tách riêng: phổ biên độ
\textbar X(f, t)\textbar{} được sử dụng làm đầu vào cho mạng neural,
trong khi phổ pha ∠X(f, t) được giữ nguyên để tái sử dụng trong quá
trình tổng hợp tín hiệu.

\textbf{3.3.3.2. Chuẩn Hóa InstantLayerNormalization}

Trước khi đưa vào khối LSTM, phổ biên độ được xử lý qua lớp chuẩn hóa
InstantLayerNormalization - một biến thể của Layer Normalization được đề
xuất bởi Luo \& Mesgarani (2018). Khác với Batch Normalization xử lý
theo chiều batch, InstantLayerNormalization thực hiện chuẩn hóa trên
chiều đặc trưng (feature dimension) của từng khung thời gian một cách
độc lập, đảm bảo mô hình hoạt động ổn định bất kể độ dài chuỗi đầu vào.

Cụ thể, với mỗi khung, lớp này tính giá trị trung bình và phương sai
trên toàn bộ các bin tần số, sau đó chuẩn hóa để đưa phân phối về trung
bình 0 và phương sai 1, rồi tái tỉ lệ bằng hai tham số học được là gamma
và beta. Cơ chế này giúp mô hình bất biến với biên độ tổng thể của tín
hiệu đầu vào, tức là hoạt động tốt bất kể âm lượng ghi âm cao hay thấp.

\begin{table}[htbp]
\centering
\caption{Cấu trúc chi tiết khối Separation Kernel -- Giai đoạn 1}
\label{tab:sep-kernel}
\begin{tabular}{>{\raggedright\arraybackslash}p{0.2\linewidth} >{\centering\arraybackslash}p{0.35\linewidth} >{\raggedright\arraybackslash}p{0.35\linewidth}}
\toprule
\rowcolor{headerblue} \textcolor{white}{\textbf{Thành phần}} & \textcolor{white}{\textbf{Thông số}} & \textcolor{white}{\textbf{Vai trò}} \\
\midrule
\rowcolor{rowgray} LSTM lớp 1 & 128 units, return\_sequences = True & Học đặc trưng tần số cục bộ theo thời gian \\
Dropout & Tỉ lệ 0,25 & Giảm overfitting giữa hai lớp LSTM \\
\rowcolor{rowgray} LSTM lớp 2 & 128 units, return\_sequences = True & Tổng hợp phụ thuộc thời gian dài hạn \\
Dense & 257 đơn vị & Ánh xạ về không gian 257 chiều tương ứng với các bin tần số \\
\rowcolor{rowgray} Sigmoid & Đầu ra trong {[}0, 1{]} & Tạo mask lọc nhiễu M₁(f, t) \\
\bottomrule
\end{tabular}
\end{table}

\textbf{3.3.3.3. Ước Lượng Mask và Tổng Hợp Tín Hiệu}

Mask M₁ được nhân trực tiếp (element-wise) với phổ biên độ gốc để ước
lượng phổ biên độ của giọng nói sạch. Các giá trị mask gần 1 giữ lại
thành phần tương ứng (giọng nói), trong khi các giá trị gần 0 triệt tiêu
thành phần đó (nhiễu). Phổ biên độ sau khi lọc được kết hợp lại với phổ
pha gốc để tái tạo phổ phức, sau đó áp dụng biến đổi Fourier ngược
(iFFT) và kỹ thuật Overlap-and-Add để thu được các khung tín hiệu thời
gian y₁(t). Các khung này trở thành đầu vào trực tiếp cho giai đoạn 2.

\textbf{3.3.4. Giai Đoạn 2 - Xử Lý trong Miền Đặc Trưng}

\textbf{3.3.4.1. Encoder Conv1D}

Các khung tín hiệu y₁(t) từ giai đoạn 1 (512 chiều, tương ứng 512 mẫu
mỗi khung) được chiếu sang không gian đặc trưng 256 chiều thông qua một
lớp tích chập 1D kiểu point-wise - tức là tích chập với kernel kích
thước 1 và stride 1. Phép chiếu tuyến tính này học một cơ sở biểu diễn
(basis transformation) riêng biệt cho từng khung thời gian, không phụ
thuộc vào các khung lân cận.

Khác biệt so với bài báo gốc: Bài báo Westhausen \& Meyer (2020) mô tả
Encoder là Conv1D với kernel kích thước 32 và stride 32, tạo ra các đặc
trưng từ các đoạn tín hiệu không chồng lấn. Tuy nhiên, cài đặt trong dự
án này sử dụng kernel kích thước 1 (point-wise convolution), hoạt động
như một phép chiếu tuyến tính thuần túy trên từng vị trí thời gian. Điều
này giúp mô hình nhỏ gọn hơn và tương thích với pipeline phân khung đã
thực hiện ở bước tiền xử lý.

\begin{table}[htbp]
\centering
\caption{Thông số bộ mã hóa Encoder -- Giai đoạn 2}
\label{tab:encoder-params}
\begin{tabular}{>{\raggedright\arraybackslash}p{0.2\linewidth} >{\centering\arraybackslash}p{0.25\linewidth} >{\raggedright\arraybackslash}p{0.4\linewidth}}
\toprule
\rowcolor{headerblue} \textcolor{white}{\textbf{Tham số}} & \textcolor{white}{\textbf{Giá trị}} & \textcolor{white}{\textbf{Ghi chú}} \\
\midrule
\rowcolor{rowgray} Số bộ lọc & 256 & Chiều không gian đặc trưng đầu ra \\
Kích thước kernel & 1 (point-wise) & Chiếu tuyến tính từng vị trí thời gian \\
\rowcolor{rowgray} Stride & 1 & Không giảm chiều thời gian \\
Hệ số bias & Không sử dụng & use\_bias = False \\
\rowcolor{rowgray} Hàm kích hoạt & Không có & Phép chiếu tuyến tính thuần túy \\
\bottomrule
\end{tabular}
\end{table}

\textbf{3.3.4.2. InstantLayerNormalization và LSTM}

Biểu diễn mã hóa (256 chiều) được chuẩn hóa qua
InstantLayerNormalization - luôn được kích hoạt ở giai đoạn 2 bất kể cài
đặt của giai đoạn 1. Sau đó, biểu diễn được xử lý qua Separation Kernel
thứ hai với cấu trúc tương tự giai đoạn 1: hai lớp LSTM (128 units mỗi
lớp) kết hợp Dropout 0,25 ở giữa. Tuy nhiên, ở giai đoạn này mạng học
các phụ thuộc thời gian trong không gian đặc trưng 256 chiều - một biểu
diễn trừu tượng hơn so với phổ biên độ STFT.

Việc đặt LSTM vào không gian đặc trưng học được có ý nghĩa quan trọng:
mô hình không chỉ lọc nhiễu ở mức phổ tần số mà còn có thể tái cấu trúc
các mẫu cấu trúc thời gian phức tạp của giọng nói, đặc biệt hữu ích khi
xử lý các loại nhiễu có cấu trúc như âm nhạc nền hay tiếng xe cộ.

\textbf{3.3.4.3. Ước Lượng Mask, Giải Mã và Overlap-and-Add}

Đầu ra LSTM được chiếu qua lớp Dense (256 đơn vị) kết hợp Sigmoid để tạo
mask M₂ trong không gian đặc trưng. Mask được nhân với biểu diễn mã hóa
để loại bỏ các thành phần nhiễu còn sót lại từ giai đoạn 1.

Biểu diễn đã lọc sau đó được giải mã về lại miền khung thời gian (512
chiều) thông qua một lớp tích chập 1D point-wise với padding nhân quả
(causal padding). Bộ giải mã trong cài đặt này sử dụng Conv1D thông
thường - không phải Conv1D Transposed - nhờ vào kiến trúc point-wise
giúp phép chiếu ngược từ 256-D về 512-D khả thi mà không cần phép tích
chập chuyển vị.

Cuối cùng, các khung đã giải mã được ghép lại thành tín hiệu liên tục
qua kỹ thuật Overlap-and-Add với bước dịch 128 mẫu, tạo ra tín hiệu
giọng nói đầu ra ŷ(t) đã được khử nhiễu hoàn toàn qua hai giai đoạn.

\textbf{3.3.5. Phân Tích Các Thành Phần Cốt Lõi}

\textbf{3.3.5.1. Cơ Chế Dual-Signal Transformation}

Đặc trưng nổi bật nhất của DTLN là sự kết hợp hai phép biến đổi tín hiệu
có bản chất khác nhau. Biến đổi STFT cung cấp biểu diễn có ý nghĩa vật
lý rõ ràng, mỗi bin tần số tương ứng với một dải tần cụ thể, giúp mạng
học cách lọc nhiễu theo phổ tần số một cách trực quan. Trong khi đó, bộ
mã hóa Conv1D học một cơ sở biểu diễn tối ưu từ dữ liệu, cho phép mô
hình nắm bắt các đặc trưng phi tuyến mà biến đổi STFT cố định không biểu
diễn được.

\begin{table}[htbp]
\centering
\caption{So sánh hai phép biến đổi tín hiệu trong DTLN}
\label{tab:dual-signal}
\begin{tabular}{>{\raggedright\arraybackslash}p{0.2\linewidth} >{\centering\arraybackslash}p{0.33\linewidth} >{\centering\arraybackslash}p{0.33\linewidth}}
\toprule
\rowcolor{headerblue} \textcolor{white}{\textbf{Tiêu chí}} & \textcolor{white}{\textbf{STFT (Giai đoạn 1)}} & \textcolor{white}{\textbf{Conv1D Encoder (Giai đoạn 2)}} \\
\midrule
\rowcolor{rowgray} Loại biến đổi & Cố định (fixed basis) & Học được từ dữ liệu \\
Chiều đầu ra & 257 bin tần số & 256 đặc trưng ẩn \\
\rowcolor{rowgray} Ý nghĩa vật lý & Rõ ràng, diễn giải được & Trừu tượng \\
Độ linh hoạt & Cố định, không thay đổi & Thích ứng theo dữ liệu huấn luyện \\
\rowcolor{rowgray} Vai trò & Loại bỏ nhiễu tổng thể ở mức phổ & Tinh chỉnh và phục hồi chi tiết giọng nói \\
\bottomrule
\end{tabular}
\end{table}


\textbf{3.3.5.2. Multiplicative Masking - Lọc Mềm}

Cả hai giai đoạn đều sử dụng phương pháp ước lượng mask nhân
(multiplicative masking) thay vì dự đoán trực tiếp phổ giọng nói sạch
(spectral mapping). Cụ thể, mạng học ước lượng một mask có giá trị trong
{[}0, 1{]} rồi nhân trực tiếp với biểu diễn đầu vào. Giá trị mask gần 1
giữ lại thành phần giọng nói, giá trị gần 0 triệt tiêu nhiễu - mô hình
hoạt động như một bộ lọc mềm (soft filter) điều khiển bởi nội dung tín
hiệu.

Phương pháp này có hai ưu điểm chính so với spectral mapping trực tiếp.
Thứ nhất, bài toán học mask bị chặn trong {[}0, 1{]} ổn định và dễ tối
ưu hóa hơn. Thứ hai, mask nhân đảm bảo năng lượng đầu ra không bao giờ
vượt quá năng lượng đầu vào - điều quan trọng để tránh các hiện vật âm
thanh (audio artifacts) trong thực tế.

\textbf{3.3.5.3. Hàm Mất Mát Negative SNR}

Mô hình được tối ưu hóa bằng hàm mất mát Negative SNR (Signal-to-Noise
Ratio), được tính theo công thức:

\textbf{\emph{SNR = E{[}s²{]} / E{[}(s − ŝ)² + ε{]}}}

\textbf{\emph{Loss = −10 · log₁₀(SNR)}}

Trong đó s là tín hiệu giọng nói sạch tham chiếu, ŝ là tín hiệu ước
lượng đầu ra của mô hình, và ε = 10⁻⁷ là hằng số nhỏ tránh chia cho 0.
Hàm mất mát tính tỉ số giữa năng lượng tín hiệu sạch và năng lượng sai
số khử nhiễu, chuyển sang đơn vị dB rồi đảo dấu để tối thiểu hóa. Giá
trị loss càng thấp (âm hơn) tương ứng với chất lượng khử nhiễu càng tốt.

So sánh với SI-SDR: SNR đơn giản hơn SI-SDR ở chỗ không có bước chiếu tỉ
lệ (scale projection). Tuy nhiên, trong bài toán khử nhiễu giọng nói khi
biên độ đầu ra cần khớp với tín hiệu sạch tham chiếu, SNR vẫn phản ánh
hiệu quả chất lượng ước lượng và phù hợp với yêu cầu huấn luyện thực tế
của dự án.

\textbf{3.3.6. Cấu Hình Huấn Luyện trên Bộ Dữ Liệu VIVOS}

\textbf{3.3.6.1. Chiến Lược Huấn Luyện}

Bộ dữ liệu huấn luyện được xây dựng bằng cách kết hợp tập âm thanh sạch
từ VIVOS với nhiễu từ DNS Challenge theo các tỉ lệ SNR ngẫu nhiên, tạo
ra các cặp (noisy, clean) song song. Mỗi mẫu được cắt thành đoạn 3 giây
- phù hợp với độ dài trung bình của các file ghi âm trong VIVOS (3 đến 5
giây) - giúp mô hình học được ngữ cảnh đủ dài trong khi vẫn kiểm soát
được bộ nhớ GPU.

\begin{table}[htbp]
\centering
\caption{Cấu hình huấn luyện mô hình DTLN trên VIVOS}
\label{tab:training-config}
\begin{tabular}{>{\raggedright\arraybackslash}p{0.25\linewidth} >{\centering\arraybackslash}p{0.3\linewidth} >{\raggedright\arraybackslash}p{0.33\linewidth}}
\toprule
\rowcolor{headerblue} \textcolor{white}{\textbf{Tham số}} & \textcolor{white}{\textbf{Giá trị}} & \textcolor{white}{\textbf{Giải thích}} \\
\midrule
\rowcolor{rowgray} Batch size & 32 & Cân bằng bộ nhớ GPU và tính ổn định gradient \\
Độ dài đoạn âm thanh & 3 giây & Phù hợp phân bố độ dài file VIVOS \\
\rowcolor{rowgray} Số epoch tối đa & 50 & Kết hợp với cơ chế dừng sớm \\
Learning rate khởi đầu & $10^{-3}$ & Giảm dần theo ReduceLROnPlateau \\
\rowcolor{rowgray} Hệ số giảm learning rate & 0,5 mỗi khi val loss không cải thiện sau 3 epoch & ReduceLROnPlateau, factor = 0,5 \\
Dừng sớm (Early Stopping) & Sau 10 epoch không cải thiện val\_loss & Tránh overfitting \\
\rowcolor{rowgray} Optimizer & Adam với clipnorm = 3,0 & Cắt gradient tránh bùng nổ \\
Random seed & 42 & Đảm bảo khả năng tái lập kết quả \\
\bottomrule
\end{tabular}
\end{table}

\textbf{3.3.6.2. Cơ Chế Lưu Checkpoint và Theo Dõi Huấn Luyện}

Quá trình huấn luyện được giám sát qua bốn cơ chế chính.

\begin{itemize}
\tightlist
\item
  Đầu tiên, một callback tùy chỉnh lưu checkpoint sau mỗi epoch và giữ
  lại mô hình tốt nhất theo validation loss.\\
\item
  Thứ hai, ReduceLROnPlateau tự động giảm learning rate 50\% khi
  validation loss ngưng cải thiện sau 3 epoch liên tiếp, giúp mô hình
  thoát khỏi các điểm cực tiểu cục bộ.\\
\item
  Thứ ba, EarlyStopping dừng huấn luyện sau 10 epoch không có cải thiện
  để tránh overfitting.\\
\item
  Cuối cùng, CSVLogger ghi lại toàn bộ lịch sử huấn luyện (loss,
  val\_loss, learning rate) theo từng epoch để phân tích sau.
\end{itemize}

\textbf{3.3.7. Pipeline Triển Khai Thực Tế}

\textbf{3.3.7.1. Xuất Mô Hình TF Lite}

Sau khi huấn luyện, mô hình được chuyển đổi sang định dạng TF Lite để
tối ưu hóa cho suy luận trên thiết bị nhúng. Do kiến trúc hai giai đoạn
có trạng thái (stateful), mô hình được xuất thành hai sub-model riêng
biệt (Core 1 tương ứng Stage 1 và Core 2 tương ứng Stage 2), mỗi
sub-model nhận và trả về trạng thái ẩn LSTM (h, c) tại mỗi bước suy
luận. Tùy chọn dynamic range quantization có thể được áp dụng để giảm
thêm kích thước mô hình.

\textbf{3.3.7.2. Suy Luận Thời Gian Thực}

Trong chế độ suy luận thời gian thực trên Raspberry Pi 4, mô hình xử lý
từng khung 512 mẫu (32 ms) một lần. Trạng thái ẩn LSTM được duy trì liên
tục giữa các lần suy luận, đảm bảo ngữ cảnh không bị gián đoạn trong
suốt luồng âm thanh streaming. Luồng xử lý tuần tự như sau:\\
microphone thu tín hiệu → phân khung 512 mẫu → STFT → Core 1 (Stage 1) →
iFFT → Core 2 (Stage 2) → Overlap-and-Add → tín hiệu sạch đầu ra.

\textbf{3.3.8. Đánh Giá Ưu Điểm và Hạn Chế}

\textbf{3.3.8.1. Ưu Điểm}

\begin{itemize}
\tightlist
\item
  Kiến trúc nhỏ gọn (\textasciitilde1,8M tham số): Đạt hiệu năng cạnh
  tranh so với các mô hình lớn hơn nhiều lần, phù hợp để triển khai trên
  thiết bị nhúng tài nguyên hạn chế.\\
\item
  Xử lý thời gian thực (\textasciitilde32 ms latency): Đáp ứng yêu cầu
  độ trễ thấp cho các ứng dụng giao tiếp tương tác.\\
\item
  Dual-domain processing: Khai thác bổ sung hai miền biểu diễn giúp khử
  nhiễu toàn diện hơn so với phương pháp đơn miền.\\
\item
  Tối ưu cho tiếng Việt: Được huấn luyện trực tiếp trên VIVOS, phù hợp
  với đặc trưng âm học của giọng nói tiếng Việt.\\
\item
  Hỗ trợ TF Lite và stateful inference: Tương thích triển khai streaming
  real-time trên nhiều nền tảng phần cứng khác nhau.
\end{itemize}

\textbf{3.3.8.1. Hạn Chế}

\begin{itemize}
\tightlist
\item
  Chỉ hỗ trợ mono channel: Chưa khai thác được thông tin không gian từ
  mảng microphone hay kỹ thuật beamforming.\\
\item
  Nhạy cảm với nhiễu phi tĩnh: Hiệu năng có thể giảm với các loại nhiễu
  biến đổi nhanh hoặc có cấu trúc âm nhạc phức tạp.\\
\item
  Yêu cầu dữ liệu song song: Cần các cặp (noisy, clean) chất lượng cao,
  khó thu thập trong điều kiện thực tế đối với tiếng Việt.\\
\item
  Pha không được tối ưu hóa: Giai đoạn 1 tái sử dụng pha gốc thay vì ước
  lượng, có thể gây hiện tượng artifact trong trường hợp nhiễu nặng.\\
\item
  Không sử dụng windowing trong STFT: Có thể gây rò rỉ phổ (spectral
  leakage), bù đắp một phần bởi overlap cao và bộ mã hóa giai đoạn 2.
\end{itemize}

\textbf{3.3.9. Kết Luận}

DTLN là kiến trúc khử nhiễu giọng nói hiệu quả, kết hợp xử lý miền tần
số thông qua STFT và miền đặc trưng học được thông qua bộ mã hóa Conv1D
point-wise. Cơ chế hai giai đoạn với LSTM và Multiplicative Masking cho
phép mô hình phân tách giọng nói khỏi nhiễu từ hai góc độ biểu diễn bổ
sung, đạt hiệu năng cạnh tranh với tổng số tham số chỉ khoảng 1,8 triệu.

Với khả năng xử lý thời gian thực (\textasciitilde32 ms latency), hỗ trợ
xuất TF Lite cho suy luận stateful, và việc đã được huấn luyện trực tiếp
trên bộ dữ liệu VIVOS tiếng Việt, DTLN là lựa chọn phù hợp cho hệ thống
hỗ trợ bán hàng thông minh nhúng triển khai trên Raspberry Pi 4, nơi yêu
cầu cân bằng chặt chẽ giữa hiệu năng khử nhiễu và tài nguyên tính toán
hạn chế.

\textbf{3.4. Kiến trúc các mô hình xử lý ngôn ngữ (PhoWhisper \&
RoBERTa)}

\begin{itemize}
\tightlist
\item
  PhoWhisper là phiên bản tinh chỉnh (fine-tuned) của mô hình Whisper,
  được thiết kế đặc biệt để tối ưu hóa khả năng nhận dạng giọng nói cho
  tiếng Việt với nhiều giọng vùng miền khác nhau.\\
\item
  Model Whisper:
\end{itemize}

\textbf{3.5. Giao thức truyền thông (Communication Protocols)}

\end{document}
